\chapter{Einstellungen}
\label{ch:einstellungen}

\section{Wo was eingestellt wird}

\begin{table}[htbp]
  \centering
  \small
  \begin{tabularx}{\textwidth}{@{}l X@{}}
    \toprule
    Ort & Inhalt \\
    \midrule
    \pfad{bc/config.py} & alle Konstanten mit Begründung im Kommentar: Takt, Tempo, Rauschen,
      Versatz, Lagen, Grenzwerte, Modell. Gilt für neue Aufnahmen und Trainings. \\
    \pfad{einstellungen/hindernisse.json} & feste Hindernisse der Zelle (Kamerastange,
      Prüfvorrichtung, Roboterbasis) mit Mindestabständen; eingecheckt \\
    \pfad{einstellungen/oberflaeche.local.json} & zuletzt gewählte Werte der Oberfläche; lokal \\
    \pfad{einstellungen/metadaten_auswahl.json} & Vorschlagslisten für die Metadaten \\
    \pfad{sequences/*.json} & Ablauf: Punkte, Bewegungsart, Überschleifen, Tischpunkte,
      Sperrzonen \\
    \pfad{checkpoints/*/policy/bc_policy.json} & Einstellungen, mit denen ein Modell aufgenommen
      wurde; die Ausführung übernimmt sie \\
    \pfad{code/config/system.yaml} & Einstellungen der Gesamt-GUI, darunter \pfad{m4.checkpoint}
      und die Kamera oben \\
    \bottomrule
  \end{tabularx}
\end{table}

\hinweis{Werte, die mitgelernt werden (Tempo, Override, Sollwertrate, Kameras), nur zwischen
Datensätzen ändern, nie mitten in einem.}

\section{Einstellungen im Reiter Aufnahme}
\label{sec:aufnahme-einst}

\begin{table}[htbp]
  \centering
  \small
  \begin{tabularx}{\textwidth}{@{}>{\raggedright\arraybackslash}p{40mm} l X@{}}
    \toprule
    Feld & Standard & Wirkung \\
    \midrule
    Tempo Transit / Endanflug / PTP & 0,15 / 0,02 m/s, 0,5 rad/s & Geschwindigkeit des Planers;
      wird mitgelernt \\
    servo\_j-Rate & 60 Hz & Rate des Sollwertstroms \\
    Rauschen bis & 1,0 & obere Grenze des Rauschfaktors je Episode \\
    Faktor nicht je Episode ziehen & aus & fester Faktor statt zufällig \\
    Fahrten je Lage & -- & Wiederholungen derselben Lage mit neuem Rauschen \\
    Lagen & -- & Anzahl automatisch umgelegter Lagen \\
    Drehung $\pm$ & \ang{45} & Drehbereich neuer Lagen um die Startstellung \\
    Abstand zum Rand & 4 cm & Abstand neuer Lagen zum Rand der Tischpunkte \\
    Objekt nach jeder Episode zurücklegen & -- & Zurücklegen ohne Aufzeichnung \\
    Am Ende an neuer Lage ablegen \dots & -- & Umlegen und \code{PICK} nachführen \\
    Erste Lage & -- & Herkunft der ersten Lage \\
    Varianten-Versatz / -Drehung & -- & nach jedem Griff leicht verschoben ablegen \\
    Probefahrt vor der ersten Episode & -- & langsame Fahrt ohne Aufzeichnung \\
    Trockenlauf ohne Greifer & -- & Test ohne Greiferbefehle; Episoden werden verworfen \\
    Nach jeder Episode bewerten & -- & Rückfrage Erfolg ja/nein \\
    Kameraschutz an kritischen Randfeldern & -- & keine Lagen, an denen die Kamera
      anstoßen könnte \\
    \bottomrule
  \end{tabularx}
\end{table}

Der gezielte Versatz im Anflug ist auf der Kommandozeile einstellbar: \pfad{--anflug-anteil}
(Standard 0,5), \pfad{--anflug-versatz 5,10} (mm), \pfad{--anflug-drehung 10,20} (Grad),
\pfad{--anflug-ausreisser}.

\section{Kameraeinstellungen}
\label{sec:kamera-einst}

Fenster \knopf{Kamera \dots} im Reiter Aufnahme:
\begin{itemize}
  \item Belichtung und Gain je Kamera einstellen, \knopf{Übernehmen} zum Ausprobieren,
    \knopf{Für Aufnahmen speichern} für alle folgenden Aufnahmen.
  \item \enquote{Wie aufgezeichnet} lädt die Werte eines vorhandenen Modells, damit die Fahrt
    mit denselben Bildern läuft.
  \item Der Rahmen im Bild zeigt den Ausschnitt, den die Policy sieht.
  \item Automatik vermeiden: Sie ändert die Bildhelligkeit zwischen Sitzungen.
\end{itemize}
Weißabgleich ist je Seriennummer fest in \pfad{config.CAMERA_WHITE_BALANCE} hinterlegt.

\hinweis{Steht der Gain der Handgelenkkamera am Anschlag (24 dB), lieber die Belichtung erhöhen,
ohne dass das Bauteil ausbrennt.}
